\hypertarget{appendix-ef4d27797a92}{}\label{referencestart:appendix-ef4d27797a92}
\section{Bits, Numbers and Memory - Visual Recall}\label{appendix-ef4d27797a92-bits-numbers-and-memory---visual-recall}

\hypertarget{appendix-ef4d27797a92-bits-numbers-and-memory---visual-recall}{}

\textbf{Scope:} Parts 1--100 only, M001--M006. Companion to {From Signals to Meaning}. No C++ Volume 02 material is included. This is a retrieval desk, not another textbook.

\subsection{How to use this desk}\label{appendix-ef4d27797a92-how-to-use-this-desk}

\hypertarget{appendix-ef4d27797a92-how-to-use-this-desk}{}

Cover the visual, sketch it from memory, answer the numbered prompts, then expand the answer key. Score each prompt \textbf{0} (missing/wrong), \textbf{1} (correct label but weak explanation), or \textbf{2} (can reconstruct an example and its limits). Revisit weak prompts in the linked teaching section. Do not turn recognition of a diagram into a mastery claim.

Mermaid shows relationships; byte strips show placement; tables show exact distinctions; state/sequence diagrams show what changes over time. Diagrams are conceptual models, not fixed CPU/OS layouts. Answer keys use collapsible Obsidian callouts.

\subsection{Short-form recall --- the whole route}\label{appendix-ef4d27797a92-short-form-recall-the-whole-route}

\hypertarget{appendix-ef4d27797a92-short-form-recall-the-whole-route}{}

\textbf{Pattern is not meaning.} A representation contract assigns weights and interpretation. → \textbf{Bits become decisions.} Masks select fields; packing needs widths, ranges, and a policy. → \textbf{Arithmetic is finite.} Choose wrap, saturation, rejection, or widening before calculation. → \textbf{Real values use a ruler.} Floating point trades spacing against range; fixed point chooses a scale. → \textbf{Operations expose error.} Rounding, cancellation, conditioning, and reproducibility matter. → \textbf{Keep the needed information.} Quantization/codebooks spend precision deliberately. → \textbf{Bytes occupy places.} Alignment, lifetime, ownership, and virtual mapping are distinct. → \textbf{Files are agreements.} Define schema, validate lengths, then decode and check integrity.

\pandocbounded{\includegraphics[keepaspectratio,alt={Diagram}]{../current/Assets/Print/bf04facd96de297a8981533eac4d37814ec29cfdaf42bb9440666efaff6faccd.pdf}}

\textbf{Caption:} Follow one value from physical signal to stored bytes and back. At every transition ask: \emph{Which bits? What meaning? What range/error policy? Where and during which lifetime? What format contract?} For the end-to-end laboratory, see {the complete journey}.

\begin{quote}
{[}!abstract{]} Compact teaching notes \textbf{Represent:} bits are patterns; interpretation supplies meaning. \hyperref[01-foundations-i-ch001-one-bit-several-bits-and-an-interpretation]{Bit-pattern example} · \hyperref[01-foundations-i-ch001-why-invert-and-add-one-works]{Two\textquotesingle s-complement derivation}

\textbf{Select and compute:} mask = select; shift = reposition; pack = preserve neighbors plus insert validated fields. Arithmetic needs an explicit overflow policy. \hyperref[01-foundations-i-ch001-fields-select-normalize-replace]{Field figure/formulas} · \hyperref[01-foundations-i-ch001-test-the-packed-meaning]{Packing lab} · \hyperref[01-foundations-i-ch001-choose-the-overflow-policy-before-calculating]{Overflow choices}

\textbf{Approximate:} floating point has nonuniform spacing; fixed point has a chosen scale. Rounding is part of each operation; tolerance is a domain rule, not a universal epsilon. \hyperref[01-foundations-i-ch001-encode-10.625-then-decode-13.25]{Worked encoding} · \hyperref[01-foundations-i-ch001-cancellation-can-expose-errors-made-earlier]{Cancellation figure} · \hyperref[01-foundations-i-ch001-a-fixed-ruler-for-fractional-values]{Fixed-point example}

\textbf{Store and exchange:} address alone does not prove valid access; alignment, lifetime, ownership, and permissions are separate. A file needs an explicit schema and checked lengths. \href{../01-Bits-Programs-and-Reliable-Systems-P02/main.pdf\#nameddest=01-foundations-i-ch002-virtual-memory-translates-page-identity}{Address translation (Bits, Programs, and Reliable Systems, Part 2, Ch. 2)} · \href{../01-Bits-Programs-and-Reliable-Systems-P02/main.pdf\#nameddest=01-foundations-i-ch002-storage-lifetime-ownership-and-placement-are-different-axes}{Memory model (Bits, Programs, and Reliable Systems, Part 2, Ch. 2)} · \href{../01-Bits-Programs-and-Reliable-Systems-P02/main.pdf\#nameddest=01-foundations-i-ch002-validate-lengths-before-performing-the-arithmetic-they-request}{Parser safety lab (Bits, Programs, and Reliable Systems, Part 2, Ch. 2)} · \href{../01-Bits-Programs-and-Reliable-Systems-P02/main.pdf\#nameddest=01-foundations-i-ch002-read-a-dump-by-reversing-the-contract}{Hex-dump lab (Bits, Programs, and Reliable Systems, Part 2, Ch. 2)}
\end{quote}

\subsection{Figure-linked pocket notes}\label{appendix-ef4d27797a92-figure-linked-pocket-notes}

\hypertarget{appendix-ef4d27797a92-figure-linked-pocket-notes}{}

Read one row, reconstruct its example, then follow the \textbf{figure} for the diagram and caption or \textbf{explanation} for the full teaching account. These four figures also live inside the continuous chapter, so the teaching text and recall desk share one reference. Each figure is a separate compact horizontal panel; the complete volume covers more than one printable page.

{\def\LTcaptype{none} % do not increment counter
\begin{longtable}[]{@{}
  >{\raggedright\arraybackslash}p{(\linewidth - 4\tabcolsep) * \real{0.3333}}
  >{\raggedright\arraybackslash}p{(\linewidth - 4\tabcolsep) * \real{0.3333}}
  >{\raggedright\arraybackslash}p{(\linewidth - 4\tabcolsep) * \real{0.3333}}@{}}
\toprule\noalign{}
\begin{minipage}[b]{\linewidth}\raggedright
Topic
\end{minipage} & \begin{minipage}[b]{\linewidth}\raggedright
Smallest useful note
\end{minipage} & \begin{minipage}[b]{\linewidth}\raggedright
Return to the teaching chapter
\end{minipage} \\
\midrule\noalign{}
\endhead
\bottomrule\noalign{}
\endlastfoot
Interpretation & Same bits, different weights: \BookInlineCode{AD} is unsigned 173 or signed −83. State width and encoding first. & \hyperref[01-foundations-i-ch001-figure-b1---one-pattern-two-integer-readings]{Figure B1} · \hyperref[01-foundations-i-ch001-the-same-bits-can-mean-different-integers]{Explanation} \\
Field replacement & Validate → clear → position → OR. \BookInlineCode{AD→9D} changes quality 2→1; outside-mask bits remain \BookInlineCode{8D}. & \hyperref[01-foundations-i-ch001-figure-b2---replace-a-field-without-damaging-its-neighbors]{Figure B2} · \hyperref[01-foundations-i-ch001-fields-select-normalize-replace]{Formulas and trace} \\
Finite arithmetic & Exact \(250+20=270\): eight-bit wrap gives 14, saturation 255, rejection no value. Widen before evaluating. & \hyperref[01-foundations-i-ch001-choose-the-overflow-policy-before-calculating]{Policy table} \\
Floating fields & Classify exponent/fraction first; normal binary32 has 24 precision bits, including its implicit leading one. & \hyperref[01-foundations-i-ch001-the-two-formats-and-their-interpretation]{Field table} · \hyperref[01-foundations-i-ch001-encode-10.625-then-decode-13.25]{Encoding trace} \\
Rounding graph & FMA rounds once; separate multiply/add rounds twice. For the stated example: \(2^{-46}\) versus zero. & \hyperref[01-foundations-i-ch001-figure-b3---rounding-points-change-the-answer]{Figure B3} · \hyperref[01-foundations-i-ch001-one-rounding-rather-than-two]{Explanation} \\
Error & Cancellation can expose previous input errors even when subtraction of stored operands is exact. Tolerance needs units and a policy. & \hyperref[01-foundations-i-ch001-cancellation-can-expose-errors-made-earlier]{Worked trace} · \hyperref[01-foundations-i-ch001-epsilon-is-one-spacing-not-a-tolerance]{Spacing table} \\
Representation choice & Fixed point chooses a constant step; floats choose a changing step; quantization chooses a finite codebook. Packing valid codes loses no further information. & \hyperref[01-foundations-i-ch001-choose-the-contract-before-choosing-the-type]{Comparison} · \hyperref[01-foundations-i-ch001-an-int4-group-from-numbers-to-nibbles-and-back]{Nibble trace} \\
Translation & Mapping changes page identity and preserves offset: \BookInlineCode{12345ABC→987ABC} under the example 4-KiB mapping. & \href{../01-Bits-Programs-and-Reliable-Systems-P02/main.pdf\#nameddest=01-foundations-i-ch002-figure-b4---translation-preserves-the-offset}{Figure B4 (Bits, Programs, and Reliable Systems, Part 2, Ch. 2)} · \href{../01-Bits-Programs-and-Reliable-Systems-P02/main.pdf\#nameddest=01-foundations-i-ch002-virtual-memory-translates-page-identity}{Explanation (Bits, Programs, and Reliable Systems, Part 2, Ch. 2)} \\
Valid access & Bounds, alignment, lifetime, permitted access and mapping are separate obligations; old bytes prove none alone. & \href{../01-Bits-Programs-and-Reliable-Systems-P02/main.pdf\#nameddest=01-foundations-i-ch002-storage-lifetime-ownership-and-placement-are-different-axes}{Lifetime/ownership table (Bits, Programs, and Reliable Systems, Part 2, Ch. 2)} \\
External bytes & Specify schema → validate arithmetic/limits → decode. Integrity, authenticity and durability require their own contracts. & \href{../01-Bits-Programs-and-Reliable-Systems-P02/main.pdf\#nameddest=01-foundations-i-ch002-validate-lengths-before-performing-the-arithmetic-they-request}{Bounds rules (Bits, Programs, and Reliable Systems, Part 2, Ch. 2)} · \href{../01-Bits-Programs-and-Reliable-Systems-P02/main.pdf\#nameddest=01-foundations-i-ch002-read-a-dump-by-reversing-the-contract}{Full dump (Bits, Programs, and Reliable Systems, Part 2, Ch. 2)} \\
\end{longtable}
}

\textbf{Recall checkpoint:} draw B1--B4 without looking. Label one assumption and one failure in each. For example: decoding \BookInlineCode{AD} as a float is the wrong contract; OR-only insertion cannot clear a field; contraction can erase the intended two-rounding path; a valid address translation does not establish a live C++ object.

\subsection{The entire story on one board --- archived original}\label{appendix-ef4d27797a92-the-entire-story-on-one-board-archived-original}

\hypertarget{appendix-ef4d27797a92-the-entire-story-on-one-board-archived-original}{}

The original vertical overview is retained here for comparison; use the horizontal map above for the one-page route.

\pandocbounded{\includegraphics[keepaspectratio,alt={Diagram}]{../current/Assets/Print/b0d43367d47bec31e4d7dca2b25bfc0f918723efe9098ce9a0a775328acec481.pdf}}

\textbf{Opening challenge:} Explain each arrow without using ``it just works.'' At every stage ask: \textbf{Which bits? Which meaning? Where? Valid to access?}

\subsection{1 - A pattern does not announce its meaning}\label{appendix-ef4d27797a92-a-pattern-does-not-announce-its-meaning}

\hypertarget{appendix-ef4d27797a92-a-pattern-does-not-announce-its-meaning}{}

\hyperref[01-foundations-i-ch001-from-a-signal-to-an-interpretation]{Repair this foundation} · \hyperref[appendix-f414bd4ed05c]{Interpretation contract} · \hyperref[appendix-4fbc03d20829]{Two\textquotesingle s complement}

\begin{verbatim}
bit index       7    6    5    4    3    2    1    0
pattern         1    0    1    0    1    1    0    1
unsigned weight 128  64   32   16    8    4    2    1
signed weight  -128  64   32   16    8    4    2    1

same pattern:   binary 10101101 = hex AD = octal 255
unsigned:       128 + 32 + 8 + 4 + 1 = 173
two's complement: -128 + 32 + 8 + 4 + 1 = -83
\end{verbatim}

{\def\LTcaptype{none} % do not increment counter
\begin{longtable}[]{@{}lrll@{}}
\toprule\noalign{}
Eight-bit interpretation & Zero patterns & Range & −5 \\
\midrule\noalign{}
\endhead
\bottomrule\noalign{}
\endlastfoot
Unsigned & 1 & 0\ldots255 & Not in domain \\
Signed magnitude & 2 & −127\ldots127 & 10000101 \\
One's complement & 2 & −127\ldots127 & 11111010 \\
Two's complement & 1 & −128\ldots127 & 11111011 \\
\end{longtable}
}

For \(N\) independent bits: \(2^N\) patterns. Unsigned maximum is \(2^N-1\), not \(2^N\). Two's-complement decode subtracts \(2^N\) when the top bit is set.

\textbf{Recall}

\begin{enumerate}
\def\labelenumi{\arabic{enumi}.}
\tightlist
\item
  Why can imperfect voltage ranges represent a stable bit? Distinguish noise margin, bit, byte, octet, word and character.
\item
  Encode 13 by repeated division; decode 1101 using Horner steps. Why do hex and octal group four and three bits?
\item
  Derive invert-plus-one rather than reciting it. Why is the signed minimum asymmetric and unsafe to negate in its own signed type?
\end{enumerate}

\begin{quote}
{[}!answer{]}- Reveal 1--3 Noise margin tolerates disturbances within a logical range. A bit has two logical states; octet means exactly eight bits; C++ byte has CHAR\_BIT bits and sizeof(char)=1. A word is architecture-specific; a character can need multiple bytes. Historical byte sizes varied.

13→6 remainder1→3 remainder0→1 remainder1→0 remainder1; reverse gives1101. Horner states1→3→6→13. Bases16 and8 are powers2\^{}4 and2\^{}3.

\(2^N-x=(2^N-1-x)+1\): width-limited complement plus one. Zero negates back to zero after the carry is discarded. The positive counterpart of −128 is128, outside the signed eight-bit range. Promotion can change the actual C++ operation type; inspect it first.
\end{quote}

\subsection{2 - Select, move, replace}\label{appendix-ef4d27797a92-select-move-replace}

\hypertarget{appendix-ef4d27797a92-select-move-replace}{}

\hyperref[01-foundations-i-ch001-bits-become-decisions-and-compact-fields]{Repair decisions and fields} · \hyperref[appendix-dbe5905163c2]{Bit mask} · \hyperref[appendix-1f9963e1a5d1]{Bit packing}

\pandocbounded{\includegraphics[keepaspectratio,alt={Diagram}]{../current/Assets/Print/605208e854e520427138c17a1aa5917b71163b5845eadf237ed86c614b45a98d.pdf}}

{\def\LTcaptype{none} % do not increment counter
\begin{longtable}[]{@{}lll@{}}
\toprule\noalign{}
Operation & Eight-bit cue & Reconstruct the purpose \\
\midrule\noalign{}
\endhead
\bottomrule\noalign{}
\endlastfoot
AND & AD \& 07 = 05 & Keep selected bits \\
OR & AD OR 40 = ED & Set selected bits \\
XOR & AD XOR 08 = A5 & Toggle; repeat restores AD \\
NOT & NOT₈(0F) = F0 & Invert at a stated width \\
Extract & (AD \textgreater\textgreater{} 4) \& 03 = 2 & Move then isolate \\
Replace & quality2→1: AD→9D & Clear then insert; preserve other bits \\
\end{longtable}
}

For width \(w\), offset \(s\), carrier \(N\): require \(0<w\), \(s+w\le N\), valid shift counts, and a declared reject/truncate policy. The full-width low-mask case needs care in executable code: shifting 1 by the carrier width is invalid.

\textbf{Recall}

\begin{enumerate}
\def\labelenumi{\arabic{enumi}.}
\setcounter{enumi}{3}
\tightlist
\item
  Prove both De Morgan identities from four input pairs. Why is built-in logical AND not interchangeable with bitwise AND?
\item
  Reconstruct extraction/insertion formulas. Explain any-selected versus all-selected flags and how to test a packer beyond a round trip.
\item
  On the C++23 baseline, why can −3/2 and −3\textgreater\textgreater1 differ? What changes when a byte operand is promoted?
\end{enumerate}

\begin{quote}
{[}!answer{]}- Reveal 4--6 NOT(A AND B)=NOT A OR NOT B; NOT(A OR B)=NOT A AND NOT B. Logical operators test truth and built-in AND/OR short-circuit; bitwise operators transform every bit.

Low mask \(L=2^w-1\), positioned mask \(M=L2^s\). Extract is right-shift then AND L; insert clears M and ORs the shifted valid field. Any tests nonzero; all compares with the whole mask. Use independent golden bytes, preservation/range checks and boundary/invalid cases: encoder and decoder can share one bug.

Integer division truncates toward zero (−1); arithmetic right shift rounds downward (−2). Promotions determine the operation width, including the width of NOT; a narrow stored type does not force a narrow expression.
\end{quote}

\subsection{3 - The adder is simple; the policy is not}\label{appendix-ef4d27797a92-the-adder-is-simple-the-policy-is-not}

\hypertarget{appendix-ef4d27797a92-the-adder-is-simple-the-policy-is-not}{}

\hyperref[01-foundations-i-ch001-arithmetic-becomes-a-policy]{Repair finite arithmetic} · \hyperref[appendix-d1a5b1e520c2]{Full adder} · \hyperref[appendix-2ed044358b82]{Arithmetic policies}

\pandocbounded{\includegraphics[keepaspectratio,alt={Diagram}]{../current/Assets/Print/c67a439285034ae2d49bf1c8993e36c97a54ea7059286994e9ebdef4fcfcaa16.pdf}}

{\def\LTcaptype{none} % do not increment counter
\begin{longtable}[]{@{}
  >{\raggedright\arraybackslash}p{(\linewidth - 4\tabcolsep) * \real{0.3000}}
  >{\raggedleft\arraybackslash}p{(\linewidth - 4\tabcolsep) * \real{0.4000}}
  >{\raggedright\arraybackslash}p{(\linewidth - 4\tabcolsep) * \real{0.3000}}@{}}
\toprule\noalign{}
\begin{minipage}[b]{\linewidth}\raggedright
Calculation
\end{minipage} & \begin{minipage}[b]{\linewidth}\raggedleft
Exact/wide result
\end{minipage} & \begin{minipage}[b]{\linewidth}\raggedright
Eight-bit destination policy
\end{minipage} \\
\midrule\noalign{}
\endhead
\bottomrule\noalign{}
\endlastfoot
250+20 & 270 & Wrap14 / saturate255 / reject \\
127+1 & 128 & Signed eight-bit overflow; no legal wrapping assumption \\
0−1 unsigned & −1 mathematically & Residue255 \\
13×11 & 13+26+104 & 143 \\
13÷3 & \(13=4(3)+1\) & Quotient4, remainder1 \\
\end{longtable}
}

\begin{verbatim}
stored type → promotion → common operation type → calculation → destination conversion
uint8 + uint8       commonly becomes int arithmetic before narrowing
signed overflow    must be prevented BEFORE evaluating the expression
\end{verbatim}

\textbf{Recall}

\begin{enumerate}
\def\labelenumi{\arabic{enumi}.}
\setcounter{enumi}{6}
\tightlist
\item
  Derive a half adder, add carry-in, and explain ripple delay versus generate/propagate carry networks. Distinguish carry, overflow, negative and zero flags.
\item
  Rebuild shift-and-add multiplication and long division. What intermediate width and invalid inputs must be handled?
\item
  Choose wrap/saturate/reject/widen for a counter, image sample and size calculation. Why can saturation break associativity?
\item
  Explain signed minimum/−1, unsigned reverse-loop failure, plain-char signedness, exact/least/fast widths, and mixed-sign comparisons.
\end{enumerate}

\begin{quote}
{[}!answer{]}- Reveal 7--10 Half-adder sum=A XOR B, carry=A AND B. Full-adder carry is true when at least two of its three inputs are true. Ripple chains carry dependencies; prefix designs combine generate/propagate information. Carry answers an unsigned question; overflow a signed-range question; N observes a result sign bit; Z observes all-zero result. Borrow conventions vary: binary6502 SBC uses A−B−(1−C), with C=1 meaning no borrow.

Set multiplier bits select shifted partial products. Division shifts bits into a remainder, conditionally subtracts, and satisfies D=Qd+R with0≤R\textless d.~Prove partial/intermediate ranges. Reject division by0 and same-type signed-min/−1.

Policies follow meaning: bounded samples may clamp, modular counters may wrap, sizes usually need checked arithmetic. Signed overflow remains undefined behavior; widening after the overflowing expression is too late. Saturation and capped updates are order-dependent.

An unsigned i≥0 never terminates a reverse loop. char signedness varies; exact-width aliases are optional, least/fast may be wider. Promotions/common-type conversions can turn a negative operand into a large unsigned value. Track all five stages, not only storage width.
\end{quote}

\subsection{4 - A stretching ruler with special regions}\label{appendix-ef4d27797a92-a-stretching-ruler-with-special-regions}

\hypertarget{appendix-ef4d27797a92-a-stretching-ruler-with-special-regions}{}

\hyperref[01-foundations-i-ch001-a-stretching-ruler-for-real-valued-quantities]{Repair floating representation} · \hyperref[appendix-343615f49ddf]{IEEE-754 floating point}

\begin{verbatim}
binary32: [ sign:1 ][ stored exponent:8 ][ fraction:23 ]  normal precision24
binary64: [ sign:1 ][ stored exponent:11][ fraction:52 ]  normal precision53
\end{verbatim}

\pandocbounded{\includegraphics[keepaspectratio,alt={Diagram}]{../current/Assets/Print/3a8ad3affccdd7a968983691ee7c2c4167d3de29ec2c42930dc1deb804b27405.pdf}}

Binary32 normal value, for \(1\le E\le254\):

\[
x=(-1)^S(1+F/2^{23})2^{E-127}.
\]

Here S is sign0/1, E stored exponent, F the unsigned23-bit fraction. Subnormal value is \((-1)^SF2^{-149}\).

{\def\LTcaptype{none} % do not increment counter
\begin{longtable}[]{@{}ll@{}}
\toprule\noalign{}
Visual checkpoint & Reconstruct \\
\midrule\noalign{}
\endhead
\bottomrule\noalign{}
\endlastfoot
10.625 & \(1.010101_2\times2^3\) → word412A0000 \\
C1540000 & Sign1, exponent130, significand1.65625 → −13.25 \\
Above1 & binary32 spacing \(2^{-23}\) \\
Immediately below1 & gap \(2^{-24}\) \\
Subnormal grid & step \(2^{-149}\), not a hidden-one format \\
\end{longtable}
}

\textbf{Recall}

\begin{enumerate}
\def\labelenumi{\arabic{enumi}.}
\setcounter{enumi}{10}
\tightlist
\item
  Decode the class before the normal formula. Why do +0 and −0 compare equal but differ as representations?
\item
  Reconstruct both words above and explain bias127/1023 and hidden precision24/53.
\item
  Distinguish epsilon, min-normal, min-subnormal and local ULP. Explain guard/round/sticky and all four rounding directions.
\end{enumerate}

\begin{quote}
{[}!answer{]}- Reveal 11--13 Reserved exponent fields alter the interpretation. Zero has E=F=0 and a separate sign; infinities have all-one E,F=0; NaNs have all-one E,F≠0. Signed-zero reciprocals can differ under IEEE behavior; signbit sees the sign even when x\textless0 does not. NaN comparisons are unordered; payload/signaling behavior needs a declared policy.

Normal exponent is E−bias and significand is1+F/2\^{}fractionBits. Subnormals remove the hidden one and keep a fixed tiny grid. A C++ float is not universally required to be binary32.

Epsilon is spacing above1, not a universal comparison tolerance. Binary32 min-normal=\(2^{-126}\); min-subnormal=\(2^{-149}\). Normal binade exponent e gives step \(2^{e-23}\). Guard/round/sticky summarize discarded information. Nearest-even picks the even retained significand at an exact tie; upward/downward/toward-zero are signed directions. Check and restore runtime rounding state.
\end{quote}

\subsection{5 - Follow the rounding points, not algebra alone}\label{appendix-ef4d27797a92-follow-the-rounding-points-not-algebra-alone}

\hypertarget{appendix-ef4d27797a92-follow-the-rounding-points-not-algebra-alone}{}

\hyperref[01-foundations-i-ch001-when-arithmetic-meets-uncertainty]{Repair numerical reasoning} · \hyperref[appendix-a2773a6a2f9d]{Cancellation and loss of significance} · \hyperref[appendix-5f1a2c3f97ff]{Numerical reproducibility}

\pandocbounded{\includegraphics[keepaspectratio,alt={Diagram}]{../current/Assets/Print/7728ef9ce14b01199dc60af382861e305eb147f624491dec5802d6191f76109e.pdf}}

{\def\LTcaptype{none} % do not increment counter
\begin{longtable}[]{@{}
  >{\raggedright\arraybackslash}p{(\linewidth - 2\tabcolsep) * \real{0.5000}}
  >{\raggedright\arraybackslash}p{(\linewidth - 2\tabcolsep) * \real{0.5000}}@{}}
\toprule\noalign{}
\begin{minipage}[b]{\linewidth}\raggedright
Witness
\end{minipage} & \begin{minipage}[b]{\linewidth}\raggedright
What it exposes
\end{minipage} \\
\midrule\noalign{}
\endhead
\bottomrule\noalign{}
\endlastfoot
0.1 & Reduced denominator10 is not a power of two: repeating binary fraction \\
binary32 \(2^{24}+1\) & Nearest-even stores \(2^{24}\): absorption \\
Nearby stored x−y & Subtraction can be exact while intended inputs remain uncertain \\
min-normal/2 & Exact subnormal is not automatically an underflow flag \\
FMA versus multiply then add & One final rounding versus two rounding points \\
Atomic floating sum & Race-free update does not fix update order \\
\end{longtable}
}

\textbf{Recall}

\begin{enumerate}
\def\labelenumi{\arabic{enumi}.}
\setcounter{enumi}{13}
\tightlist
\item
  When does a rational terminate in binary? Why can cancellation magnify relative input error even when subtraction itself is exact?
\item
  Classify overflow, tiny exact results, tiny inexact results, invalid operations and division-by-zero flags.
\item
  Distinguish absolute/relative and forward/backward error; stability versus conditioning. Why is approximate closeness not equality?
\item
  Draw factored/distributed and FMA/unfused evaluation graphs. What do pairwise/compensated sums, wider accumulation and fixed order solve---and not solve?
\end{enumerate}

\begin{quote}
{[}!answer{]}- Reveal 14--17 A reduced rational terminates iff its denominator is a power of two. Sterbenz's nearby-float condition x/2≤y≤2x for nonnegative stored values can give exact subtraction under the stated gradual-underflow conditions. A small difference can still expose earlier input uncertainty; sensitivity scales with (\textbar x\textbar+\textbar y\textbar)/\textbar x−y\textbar.

Tiny and underflow are not synonyms. Flags also include inexact, overflow, divide-by-zero and invalid; classify the actual operation/environment.

Absolute error has units; relative error needs a nonzero reference. Forward error measures output discrepancy; backward error asks how perturbed inputs explain the output. Conditioning belongs to the problem, stability to the algorithm. Validate tolerances and non-finite inputs first; pairwise closeness need not be transitive.

FMA rounds ab+c once; unfused arithmetic rounds the product too. Reassociation/distribution move rounding points. Pairwise and compensated algorithms can reduce error but do not reconstruct lost input bits or guarantee identical cross-platform execution. Wider transient arithmetic does not preserve a repeatedly narrowed state. Determinism needs order, precision, compiler/environment and algorithm contracts; capped TSDF-style weighted updates can remain history-dependent mathematically.
\end{quote}

\subsection{6 - Choose what information to keep}\label{appendix-ef4d27797a92-choose-what-information-to-keep}

\hypertarget{appendix-ef4d27797a92-choose-what-information-to-keep}{}

\hyperref[01-foundations-i-ch001-choosing-what-information-to-keep]{Repair representation choices} · \hyperref[appendix-1d0dbe8924cc]{Fixed-point representation} · \hyperref[appendix-d5da3eaabaef]{Affine quantization and zero point}

{\def\LTcaptype{none} % do not increment counter
\begin{longtable}[]{@{}
  >{\raggedright\arraybackslash}p{(\linewidth - 4\tabcolsep) * \real{0.3333}}
  >{\raggedright\arraybackslash}p{(\linewidth - 4\tabcolsep) * \real{0.3333}}
  >{\raggedright\arraybackslash}p{(\linewidth - 4\tabcolsep) * \real{0.3333}}@{}}
\toprule\noalign{}
\begin{minipage}[b]{\linewidth}\raggedright
Format
\end{minipage} & \begin{minipage}[b]{\linewidth}\raggedright
What it buys
\end{minipage} & \begin{minipage}[b]{\linewidth}\raggedright
What remains
\end{minipage} \\
\midrule\noalign{}
\endhead
\bottomrule\noalign{}
\endlastfoot
Fixed point I/S & Constant absolute step1/S & Bounded range; rescale/overflow policy \\
FP16 fields1/5/10 & 11-bit normal precision & Maximum65504; narrow range \\
BF16 fields1/8/7 & Wide exponent range,8-bit precision & Coarse spacing; NaN/rounding policy \\
Groupwise INT4 & Small local codebook & Clipping, scales, padding, task error \\
UNORM8 & 256 codes in0\ldots1 & Step1/255; 128 is not exactly0.5 \\
SNORM8 & Normalized signed range & −128 and−127 both decode−1 \\
\end{longtable}
}

\begin{verbatim}
fixed F=8:    1.5 → raw384     2.25 → raw576
multiply:    384×576 =221184  (scale256²)
rescale:     221184/256=864   (scale256)
decode:      864/256=3.375

offset INT4 codes:  [1,5,8,11,15]
low nibble first:  [5|1] [11|8] [padding0|15] → 51 B8 0F
quantization loses information; packing valid codes does not.
\end{verbatim}

\textbf{Recall}

\begin{enumerate}
\def\labelenumi{\arabic{enumi}.}
\setcounter{enumi}{17}
\tightlist
\item
  Choose F for a signed32-bit fixed domain ±10\^{}6 with step≤0.001. Derive multiply/divide rescaling and nearest unclipped error.
\item
  Encode1.5 in FP16; compare FP16/BF16 near1 and at large magnitude. Why can truncating a NaN payload accidentally produce infinity?
\item
  Explain affine q/z/scale, outlier effects, group metadata, learned one-bit centroids, RMSE/MAE and normalized graphics/color-space distinctions.
\end{enumerate}

\begin{quote}
{[}!answer{]}- Reveal 18--20 Require \(2^{-F}\le0.001\) and \(10^6 2^F\le2^{31}-1\): F=10 or11. Multiply raw operands at scale S² then round AB/S; quotient rounds AS/B for B≠0. Prove intermediate capacity before calculating, check destination range, and define negative ties. Nearest unclipped error≤1/(2S); relative error can explode nearzero.

FP16 1.5→3E00; bias15, normal min2\^{}-14, subnormal min2\^{}-24. BF16 bias127, normal spacing above1 is2\^{}-7 versus FP16's2\^{}-10. Dropping a nonzero discarded NaN fraction can leave all-one exponent and zero fraction: infinity. Classify first; canonicalization is a policy.

Affine reconstruction is scale×(q−z), with positive valid scale; nearest unclipped ideal error≤scale/2. A max-absolute outlier can erase small values. INT4 payload costs ceil(N/2)+4ceil(N/G) bytes with float32 scales, not four bits/value universally. At N=5,G=32 it costs7 bytes. One-bit labels can select learned means rather than±1. RMSE squares errors; MAE takes magnitudes; neither proves application quality. UNORM is not sRGB; quantized normals may need safe renormalization and still lose direction.
\end{quote}

\subsection{7 - Placement, translation and lifetime are separate}\label{appendix-ef4d27797a92-placement-translation-and-lifetime-are-separate}

\hypertarget{appendix-ef4d27797a92-placement-translation-and-lifetime-are-separate}{}

\href{../01-Bits-Programs-and-Reliable-Systems-P02/main.pdf\#nameddest=01-foundations-i-ch002-where-the-bytes-live}{Repair memory reasoning (Bits, Programs, and Reliable Systems, Part 2, Ch. 2)} · \href{../01-Bits-Programs-and-Reliable-Systems-P02/main.pdf\#nameddest=appendix-bc07c3d25b18}{Endianness (Bits, Programs, and Reliable Systems, Part 2, Reference notes: Endianness)} · \href{../01-Bits-Programs-and-Reliable-Systems-P02/main.pdf\#nameddest=appendix-a1f7c53b4742}{Alignment and padding (Bits, Programs, and Reliable Systems, Part 2, Reference notes: Alignment and padding)} · \href{../01-Bits-Programs-and-Reliable-Systems-P02/main.pdf\#nameddest=appendix-01e5905165c4}{Page tables and MMU (Bits, Programs, and Reliable Systems, Part 2, Reference notes: Page tables and MMU)}

\begin{verbatim}
value 12345678 hex     increasing byte addresses →
little endian         78 56 34 12
big endian            12 34 56 78

ordinary ABI example: char / uint32 / char, alignments1/4/1
offset 0: [a][pad][pad][pad][b b b b][c][tail][tail][tail]
member offsets0,4,8; size12; alignment4; array stride12
\end{verbatim}

\pandocbounded{\includegraphics[keepaspectratio,alt={Diagram}]{../current/Assets/Print/e8467ed42e7c2ac1b5ce1bcb0132b6bd95463b0207a4de2f069d991f2d49b28f.pdf}}

Example4-KiB pages: VA12345ABC → VPN12345, offsetABC → PFN987 → PA987ABC. A TLB miss is not automatically a fault; a fault is not automatically disk I/O or a crash.

\pandocbounded{\includegraphics[keepaspectratio,alt={Diagram}]{../current/Assets/Print/533250314ebfe5d22aebb39a1cf72953db6b28af0074507e379f332ade24d194.pdf}}

\textbf{Recall}

\begin{enumerate}
\def\labelenumi{\arabic{enumi}.}
\setcounter{enumi}{20}
\tightlist
\item
  Why does swapping bytes not reverse bits? Derive aligned offsets and tail padding; compare packing, reordering, AoS/SoA and CPU/GPU layouts.
\item
  Reconstruct the address translation. Distinguish process isolation, shared mappings, TLB/data-cache misses, faults, COW, virtual/physical contiguity and huge-page tradeoffs.
\item
  Separate stack/heap/static regions from storage duration and object lifetime. Explain fragmentation, dangling views, mapped files, const, W\^{}X and ASLR.
\end{enumerate}

\begin{quote}
{[}!answer{]}- Reveal 21--23 Endianness orders significance bytes, not bits within an octet. Align each member under stated ABI requirements; round total extent for aligned array stride. Reordering can reduce holes but change ABI. Packing may misalign; smaller is not automatically faster. AoS groups records, SoA fields; hot access patterns and backend layout contracts decide.

For page size P, VPN=floor(VA/P), offset=VA mod P; PA=PFN×P+offset. Context/permissions matter. TLB caches translations; data cache caches contents. COW copies a shared backing page when a permitted write needs a private copy. Consecutive virtual pages may use scattered frames. Larger pages increase simplified TLB reach entries×pageSize but reduce granularity/flexibility.

Automatic objects often use stack frames, but optimization/ABI can differ; heap allocation manages variable-lifetime blocks. Static duration can still need runtime initialization. Internal fragmentation is allocated slack; external fragmentation is separated holes. A span does not own/extend backing life. A mapped file needs bounds, permissions and access/lifetime validity. const does not force a read-only segment; W\^{}X concerns mapping permissions; ASLR varies locations. Old bytes do not prove a live object.
\end{quote}

\subsection{8 - External bytes need a complete agreement}\label{appendix-ef4d27797a92-external-bytes-need-a-complete-agreement}

\hypertarget{appendix-ef4d27797a92-external-bytes-need-a-complete-agreement}{}

\href{../01-Bits-Programs-and-Reliable-Systems-P02/main.pdf\#nameddest=01-foundations-i-ch002-bytes-become-a-durable-agreement}{Repair external bytes (Bits, Programs, and Reliable Systems, Part 2, Ch. 2)} · \href{../01-Bits-Programs-and-Reliable-Systems-P02/main.pdf\#nameddest=appendix-652d81568de9}{Checked binary parsing (Bits, Programs, and Reliable Systems, Part 2, Reference notes: Checked binary parsing)} · \href{../01-Bits-Programs-and-Reliable-Systems-P02/main.pdf\#nameddest=appendix-cbe594d1ab0a}{Integrity and authenticity (Bits, Programs, and Reliable Systems, Part 2, Reference notes: Integrity and authenticity)}

\pandocbounded{\includegraphics[keepaspectratio,alt={Diagram}]{../current/Assets/Print/24a651e6afd06573e2dc085e569cd96adbbcc06e92825a4437d9d288b61d7ee4.pdf}}

{\def\LTcaptype{none} % do not increment counter
\begin{longtable}[]{@{}ll@{}}
\toprule\noalign{}
Same binary32-compatible value1 & Result \\
\midrule\noalign{}
\endhead
\bottomrule\noalign{}
\endlastfoot
Numeric conversion to uint32 & 1 \\
Representation bit-cast & Word3F800000 \\
Little-endian encoding & Bytes00 00 80 3F \\
Big-endian encoding & Bytes3F 80 00 00 \\
\end{longtable}
}

{\def\LTcaptype{none} % do not increment counter
\begin{longtable}[]{@{}
  >{\raggedright\arraybackslash}p{(\linewidth - 4\tabcolsep) * \real{0.3333}}
  >{\raggedright\arraybackslash}p{(\linewidth - 4\tabcolsep) * \real{0.3333}}
  >{\raggedright\arraybackslash}p{(\linewidth - 4\tabcolsep) * \real{0.3333}}@{}}
\toprule\noalign{}
\begin{minipage}[b]{\linewidth}\raggedright
Defense
\end{minipage} & \begin{minipage}[b]{\linewidth}\raggedright
Answers
\end{minipage} & \begin{minipage}[b]{\linewidth}\raggedright
Does not establish
\end{minipage} \\
\midrule\noalign{}
\endhead
\bottomrule\noalign{}
\endlastfoot
Magic/version & Which proposed format/revision? & Valid dimensions or trusted sender \\
Checked schema & Structurally permitted content? & Unchanged payload \\
CRC/checksum & Matches stated error-detection reference? & Adversarial authenticity \\
Cryptographic hash & Matches trusted digest? & Trust if attacker replaces digest too \\
MAC/signature & Authenticated under trusted key policy? & Semantic correctness \\
Atomic publication & Replacement visibility contract? & Crash durability \\
\end{longtable}
}

\textbf{Recall}

\begin{enumerate}
\def\labelenumi{\arabic{enumi}.}
\setcounter{enumi}{23}
\tightlist
\item
  State bit-cast preconditions; distinguish native ABI from external schema. Give safe multiplication and range checks before parsing.
\item
  Explain checksum blind spots, specified CRC parameters, hash collisions, authenticity and durability. Decode the mini-dump below.
\end{enumerate}

\begin{quote}
{[}!answer{]}- Reveal 24--25 Bit-cast needs same size, trivially copyable types and valid initialized representations under the applicable rules. It is not numeric conversion, endian conversion or permission for arbitrary typed access.

Before unsigned a×b check a=0 or b≤MAX/a. Before reading n bytes at o in length L, check o≤L then n≤L−o. Check header≤length, exact payload length, size\_t representability and practical resource limits before allocation.

Additive byte sums miss reorderings. A CRC needs polynomial, initial/final operations, reflection/order and covered range. CRC-32/ISO-HDLC check string123456789 gives CBF43926. Finite digests necessarily have collisions over unrestricted messages. A hostile replacement requires trusted digest/key policy; valid decoding, integrity, authenticity, publication and durability are different claims.
\end{quote}

\subsubsection{Mini-dump - Reconstruct before revealing}\label{appendix-ef4d27797a92-mini-dump---reconstruct-before-revealing}

\hypertarget{appendix-ef4d27797a92-mini-dump---reconstruct-before-revealing}{}

\begin{verbatim}
0000  4E 51 54 4E 53 52 30 31   magic
0008  01 00 00 00               version
000C  02 00 00 00 00 00 00 00   rows
0014  03 00 00 00 00 00 00 00   columns
001C  00 00 80 3F               first payload word
\end{verbatim}

Assume explicit little-endian uint32 version, uint64 dimensions and IEEE binary32 payload. What are shape, header size, payload size and total size? What changes if payload byte80 becomes00?

\begin{quote}
{[}!answer{]}- Reveal the dump Shape2×3; header28 bytes; payload24 bytes; total52 bytes. First word3F800000 decodes1. Flipping80→00 gives3F000000, value0.5. Magic, dimensions and file length can remain plausible; structural validation alone misses this substitution. This is a teaching schema, not a freshly verified production format.
\end{quote}

\subsection{Final blank-board test}\label{appendix-ef4d27797a92-final-blank-board-test}

\hypertarget{appendix-ef4d27797a92-final-blank-board-test}{}

Without looking above, draw: interpretation fork → mask/extract/replace → adder with two flag questions → floating classifier → rounding graph → precision-choice table → byte/alignment strip → virtual translation → live-object state → safe reader.

For every diagram supply \textbf{one numeric example, one invariant, one failure case and one assumption}. Then explain the same stored bytes under two interpretations. A round trip is only one test: demand independent golden bytes and invalid-input checks too.

\subsection{Coverage map - No source parts dropped}\label{appendix-ef4d27797a92-coverage-map---no-source-parts-dropped}

\hypertarget{appendix-ef4d27797a92-coverage-map---no-source-parts-dropped}{}

This map is navigation, not partwise lesson segmentation. Overlapping examples may revisit a concept; the listed coverage ranges assign each of1--100 exactly once.

{\def\LTcaptype{none} % do not increment counter
\begin{longtable}[]{@{}ll@{}}
\toprule\noalign{}
Source parts & Recall panel \\
\midrule\noalign{}
\endhead
\bottomrule\noalign{}
\endlastfoot
1--7; 21--27 & \hyperref[appendix-ef4d27797a92-a-pattern-does-not-announce-its-meaning]{\#1 - A pattern does not announce its meaning} \\
8--20 & \hyperref[appendix-ef4d27797a92-select-move-replace]{\#2 - Select, move, replace} \\
28--40 & \hyperref[appendix-ef4d27797a92-the-adder-is-simple-the-policy-is-not]{\#3 - The adder is simple; the policy is not} \\
61--76 & \hyperref[appendix-ef4d27797a92-a-stretching-ruler-with-special-regions]{\#4 - A stretching ruler with special regions} \\
77--88 & \hyperref[appendix-ef4d27797a92-follow-the-rounding-points-not-algebra-alone]{\#5 - Follow the rounding points, not algebra alone} \\
89--96 & \hyperref[appendix-ef4d27797a92-choose-what-information-to-keep]{\#6 - Choose what information to keep} \\
41--60 & \hyperref[appendix-ef4d27797a92-placement-translation-and-lifetime-are-separate]{\#7 - Placement, translation and lifetime are separate} \\
97--100 & \hyperref[appendix-ef4d27797a92-external-bytes-need-a-complete-agreement]{\#8 - External bytes need a complete agreement} \\
\end{longtable}
}

Repair exact definitions in \href{../01-Bits-Programs-and-Reliable-Systems-P03/main.pdf\#nameddest=appendix-fd090d982a15}{Supplementary/Foundations (Bits, Programs, and Reliable Systems, Part 3, Reference notes: Foundations)}; longer proofs in \href{../01-Bits-Programs-and-Reliable-Systems-P03/main.pdf\#nameddest=01-foundations-i-ch010-reconstruct-the-representation-and-lifetime-proofs}{Continuous Notes/Computing Foundations - From Bits to C++ and Systems APIs\#Reconstruct the representation and lifetime proofs (Bits, Programs, and Reliable Systems, Part 3, Ch. 8)}; extended experiments in \hyperref[appendix-4e2236f53445]{Supplementary/Code Snippets} and \hyperref[appendix-92ff77d4dfdc]{Supplementary/Worked Traces}. The full teaching explanations remain in Volume01. Future recall companions should reuse this one-per-volume structure; none is created for C++ in this request.
